\documentclass{ximera}

\newcommand{\inbetween}{\fontsize{7pt}{9pt}\selectfont}


\ifdefined\HCode
  \newcommand{\alttext}[1]{\HCode{<span class="sr-only">#1</span>}}
\else
  \newcommand{\alttext}[1]{} % Fallback for PDF view
\fi

% MATH -----------------------------------------------------------
\newcommand{\nats}{\mathbb N}
\newcommand{\ints}{\mathbb Z}
\newcommand{\rats}{\mathbb Q}
\newcommand{\reals}{\mathbb R}
\newcommand{\complex}{\mathbb C}
\newcommand{\powerset}{\mathscr P}

%------------------------------------------------------------



\newcommand{\Laplace}[1]{\ensuremath{\mathcal{L}{\left[#1\right]}}}
\newcommand{\InvLap}[1]{\ensuremath{\mathcal{L}^{-1}{\left[#1\right]}}}


\pgfplotsset{compat=1.18}

\begin{document}
\begin{abstract}
 \end{abstract}

    \title{Lab 6 -  Laplace Transforms }
    
    \maketitle

Let $f(t)$ be a piecewise continuous function on $[0,\infty)$.  The \textbf{Laplace transform} of $f$ is the function $F(s)=\displaystyle \int_0^{\infty} f(t)e^{-st}\,dt$, on the set of $s$ values for which the integral converges.

\begin{enumerate}

\item By direct calculation, determine the Laplace transform of the function graphed below (assume the function continues to be equal to zero to the right of $x=4$) for $s>0$.  Put the answer in the box provided.

\hfill{}\fbox{ \begin{minipage}{3in}\vspace{0.5in} \hfill\vspace{0.5in} \end{minipage} }

% 1-U(t-2) has a Laplace transform of 1/s-(1/s)e^(-2s) for s>0

\begin{figure}[h]
\centering
\alttext{The graph of a piecewise continuous function.}
\begin{tikzpicture}
\begin{axis}[axis lines=middle, xtick={-2,-1,...,4}, ytick={-2,-1,...,4},
xmin=-2, xmax=4, ymin=-2, ymax=4]

\addplot[blue,ultra thick,domain=0:2] {1};
\addplot[blue,ultra thick,domain=2:4] {0};
\end{axis}
\end{tikzpicture}
\caption{Graph of a piecewise continuous function}
\end{figure}



\newpage


\item  Fill-in the missing details (after each of the ellipses) to show by direct calculation that the Laplace transform of $f(t)=\cos{(\omega t)}$ where $\omega>0$, is $\displaystyle \frac{s}{s^2+\omega^2}$ for $s>0$.

If $s=0$ then $\displaystyle \int_0^{\infty} \cos{(\omega t)}e^{-st}\,dt=\int_0^{\infty} \cos{(\omega t)}\,dt$, which ... \hfill .  If\\ 

 $s<0$ then $\displaystyle \int_0^{\infty} \cos{(\omega t)}e^{-st}\,dt$ also ...\hfill.  For $s> 0$, let \\ 

$\displaystyle A=\int_0^{\infty} \cos{(\omega t)}e^{-st}\,dt$.  By ... \hfill  with $u=\cos{(\omega t)}$ \\ 

and ... \hfill we get that\\ 


 $A=$ ...  \hspace{5cm} $\displaystyle \left.\:^{\;}\right|_0^{\infty} - \frac{\omega}{s} \int_0^{\infty}$ ...\hspace{5cm} $dt$.\\ 

For $s>0$, $\displaystyle \lim_{t\rightarrow \infty} -\frac{\cos{(\omega t)}}{s}e^{-st}=$ ... \hspace{5cm}  \\ 

so $\displaystyle A=\frac{1}{s} \; -$ ...\hspace{9cm} . \\ 

 By ...\hfill , with $\displaystyle u=\sin{(\omega t)}$ and ...\hfill we get \\ 

$\displaystyle A=\frac{1}{s} \; - \; \frac{\omega}{s} [$ ...\hspace{10cm} $]$. \\ 

For $s>0$, $\displaystyle \lim_{t\rightarrow \infty} -\frac{\sin{(\omega t)}}{s}e^{-st}=$ ...\hspace{5cm}  .\\ 

Then $A=\frac{1}{s} \; - [$ ...\hspace{3cm} $]A$.  Solving this for $A$ \\

we get $A=$ \hfill,  Hence the ...\hfill of $\cos (\omega t)$ is \\ \\

 $A=$ ... \hfill for $s>0$. 
 
\newpage

\item  One can also show that $\displaystyle \Laplace{\sin (\omega t))}=\frac{\omega}{s^2+\omega^2}$ for $s>0$.    Use the First Shifting Property to derive expressions for the Laplace transforms $\Laplace{e^{at}\cos{(\omega t)}}$ and $\Laplace{e^{at}\sin{(\omega t)}}$ for $s>a$, where $a,\omega>0$ are constants.  Put the answers in the box provided.

\hfill{}\fbox{ \begin{minipage}{5in}\vspace{0.75in} \hfill\vspace{0.75in} \end{minipage} }

\newpage


\item Using $\Laplace{y^{\,\prime}}=s\Laplace{y}-y(0)$, find the Laplace transform of the solution to the IVP $y^{\,\prime}+y=e^{-t}\sin{(2t)}$, $y(0)=1$ by taking the Laplace transform of both sides of the differential equation and then solving for $\Laplace{y}$ algebraically.  Put the answer in the box provided.

\hfill{}\fbox{ \begin{minipage}{3in}\vspace{0.25in} \hfill\vspace{0.25in} \end{minipage} }





\end{enumerate}

\end{document}